\documentclass[11pt,a4paper]{article}
\usepackage[italian]{babel}
\usepackage[margin=2.5cm]{geometry}
\usepackage{parskip}
\usepackage{amsmath, amssymb, amsthm}
\usepackage{booktabs}
\usepackage{array}
\usepackage{xcolor}
\usepackage{needspace}
\usepackage{enumitem}
\usepackage{listings}
\usepackage{tikz}
\usepackage[most]{tcolorbox}
\usepackage[hidelinks]{hyperref}
\usetikzlibrary{decorations.pathreplacing, shapes.geometric, arrows.meta, positioning}

% Niente vedove/orfane: i paragrafi non lasciano righe isolate a fine/inizio pagina
\widowpenalty=10000
\clubpenalty=10000

% Elenchi più compatti
\setlist{itemsep=2pt, parsep=0pt, topsep=4pt, partopsep=0pt}

% --- Riquadri con bordo nero, senza numero (si spezzano tra due pagine) ---
% Uso: \begin{definizione} ... \end{definizione}
%      \begin{teorema}[Nome] ... \end{teorema}  (il nome è facoltativo)
\tcbset{nota/.style={enhanced jigsaw, breakable, colback=white,
  colframe=black, boxrule=0.5pt, sharp corners}}
\NewTColorBox{definizione}{}{nota,
  before upper={\textbf{Definizione.}\ }}
\NewTColorBox{teorema}{o}{nota,
  before upper={\textbf{Teorema\IfValueT{#1}{ (#1)}.}\ }}
\NewTColorBox{proposizione}{o}{nota,
  before upper={\textbf{Proposizione\IfValueT{#1}{ (#1)}.}\ }}
\NewTColorBox{proprieta}{o}{nota,
  before upper={\textbf{Proprietà\IfValueT{#1}{ (#1)}.}\ }}
\NewTColorBox{assioma}{o}{nota, boxrule=1pt,
  before upper={\textbf{Assioma\IfValueT{#1}{ (#1)}.}\ }}

% --- Ambienti semplici (senza riquadro) ---
% Il titolo sta in cima al blocco, anche se il contenuto inizia con un riquadro o
% con codice; e non resta mai da solo in fondo a una pagina.
% Uso: \begin{esempio}[Titolo facoltativo] ... \end{esempio}
\newcommand{\newrunin}[2]{%
  \NewDocumentEnvironment{#1}{o}{%
    \par\Needspace{4\baselineskip}\addvspace{\medskipamount}\noindent
    \textbf{#2}\IfValueT{##1}{ (##1)}\textbf{.}\ \ignorespaces}%
  {\par\addvspace{\medskipamount}}}
\newrunin{esempio}{Esempio}
\newrunin{esercizio}{Esercizio}
\NewTColorBox{osservazione}{}{nota,
  before upper={\textbf{Osservazione.}\ }}

% V e F nelle tavole di verità
\newcommand{\V}{\textsf{V}}
\newcommand{\F}{\textsf{F}}

% --- Codice C ---
% Uso: \begin{codice} ... \end{codice}      codice C numerato
%      \begin{comando} ... \end{comando}    comandi da terminale
%      \begin{risultato} ... \end{risultato} output di un programma
%      \lstinline|printf("ciao");|           codice dentro al testo
\lstdefinestyle{c}{
  language=C,
  basicstyle=\ttfamily\small,
  keywordstyle=\color{blue}\bfseries,
  commentstyle=\color{gray}\itshape,
  stringstyle=\color{red!70!black},
  numbers=left,
  numberstyle=\tiny\color{gray},
  numbersep=8pt,
  columns=flexible,
  keepspaces=true,
  breaklines=true,
  showstringspaces=false,
  tabsize=4,
  morekeywords={include, define},
  % lettere accentate nei commenti
  literate={à}{{\`a}}1 {è}{{\`e}}1 {é}{{\'e}}1 {ì}{{\`i}}1
           {ò}{{\`o}}1 {ù}{{\`u}}1 {È}{{\`E}}1 {À}{{\`A}}1
}
\lstset{style=c}
\newtcblisting{codice}{listing only, nota, left=6mm,
  listing options={style=c}}
\newtcblisting{comando}{listing only, nota,
  listing options={basicstyle=\ttfamily\small, numbers=none, language={},
    columns=flexible, keepspaces=true, breaklines=true}}
\newtcblisting{risultato}{listing only, enhanced jigsaw, breakable,
  colback=black!5, colframe=black, boxrule=0.5pt, sharp corners,
  listing options={basicstyle=\ttfamily\small, numbers=none, language={},
    columns=flexible, keepspaces=true, breaklines=true}}

% --- Diagrammi di flusso (simboli standard) ---
\tikzset{
  terminale/.style = {ellipse, draw, minimum width=2.5cm, minimum height=0.9cm},
  io/.style        = {trapezium, trapezium left angle=70, trapezium right angle=110,
                      draw, minimum width=2.5cm, minimum height=0.9cm},
  processo/.style  = {rectangle, draw, minimum width=2.5cm, minimum height=0.9cm},
  decisione/.style = {diamond, draw, aspect=2, minimum width=2.5cm},
  freccia/.style   = {thick, -{Stealth}}
}

\title{Porte logiche e circuiti logici}
\author{Marco Cerbella}
\date{Lezione 3}

\begin{document}
\maketitle

\section{Porte logiche}

\begin{definizione}[Porta logica]
Una \emph{porta logica} è un componente elettronico di base che opera su uno o più segnali di ingresso e genera un unico segnale in uscita. È la realizzazione fisica delle funzioni elementari introdotte dall'algebra (AND, OR, NOT).
\end{definizione}

Oltre alle porte AND, OR e NOT, si usano le porte NAND, NOR, XOR e XNOR.

\begin{center}
\begin{tabular}{cc|cccccc}
\toprule
$x_0$ & $x_1$ & AND & OR & NAND & NOR & XOR & XNOR \\
\midrule
0 & 0 & 0 & 0 & 1 & 1 & 0 & 1 \\
0 & 1 & 0 & 1 & 1 & 0 & 1 & 0 \\
1 & 0 & 0 & 1 & 1 & 0 & 1 & 0 \\
1 & 1 & 1 & 1 & 0 & 0 & 0 & 1 \\
\bottomrule
\end{tabular}
\end{center}

La porta NOT ha un solo ingresso e restituisce $x'$.

\subsection{NAND e NOR: operatori funzionalmente completi}

La porta NAND realizza $(x_0 \cdot x_1)'$ e la porta NOR realizza $(x_0 + x_1)'$.

\begin{proposizione}
NAND e NOR sono \emph{funzionalmente completi}: qualsiasi funzione logica si può realizzare usando esclusivamente porte NAND, oppure esclusivamente porte NOR, invece di AND, OR e NOT.
\end{proposizione}

\begin{esempio}[NOT con una NAND]
Collegando i due ingressi di una NAND alla stessa variabile $x$ si ottiene, per idempotenza,
\[
(x \cdot x)' = (x)' = x'.
\]
Quindi la porta NOT si realizza con una NAND.
\end{esempio}

\subsection{XOR e XNOR}

Le porte XOR e XNOR corrispondono agli operatori booleani XOR e XNOR (OR e NOR \emph{esclusivi}): realizzano rispettivamente la funzione di \emph{disuguaglianza} e di \emph{uguaglianza} tra due bit. Lo XOR vale $1$ se e solo se i due ingressi sono diversi, lo XNOR vale $1$ se e solo se sono uguali.

\subsection{Porte a $n$ ingressi}

Accanto alle porte elementari a due ingressi, in genere si assume di avere a disposizione le stesse porte con un numero arbitrario di ingressi. In una AND a tre ingressi $x_0, x_1, x_2$ l'uscita $y$ vale $1$ se e solo se tutte le variabili in ingresso valgono $1$. In una OR a quattro ingressi $x_0, \dots, x_3$ l'uscita $y$ vale $1$ se almeno una variabile in ingresso vale $1$.

\section{Circuiti logici}

Le porte elementari si collegano per realizzare funzioni più complesse. A partire da un'espressione logica $e$ (ricavata, per esempio, dalla tabella di verità di una funzione booleana) si costruisce il circuito che realizza la funzione descritta da $e$.

\begin{esempio}
Il circuito per $x + y'z$ usa una porta NOT su $y$, una AND tra $y'$ e $z$ (uscita $y'z$) e una OR tra $x$ e $y'z$.
\end{esempio}

\subsection{Ritardi e cammino critico}

Ogni porta logica ha un \emph{ritardo}: il tempo tra l'istante in cui si applicano i segnali in ingresso e l'istante in cui il segnale di uscita si stabilizza a un valore costante. Il ritardo dipende dalla realtà fisica dei componenti (i transistor).

\begin{definizione}[Cammino critico]
Il \emph{cammino critico} (\emph{critical path}) è il cammino dagli ingressi alle uscite che attraversa il maggior numero di porte in un circuito, trascurando i negatori, perché possono essere accorpati alla porta precedente o successiva.
\end{definizione}

\begin{esempio}
In una rete divisa in tre livelli di porte, i cammini critici attraversano tre porte.
\end{esempio}

\subsection{Circuiti da forme canoniche}

Le funzioni espresse in prima o in seconda forma canonica corrispondono sempre a circuiti logici \emph{a due livelli}, che usano porte a $n$ ingressi.
\begin{itemize}
    \item In \emph{prima forma canonica} il primo livello contiene le porte AND che calcolano i mintermini e il secondo livello un'unica porta OR che ne fa la somma.
    \item In \emph{seconda forma canonica} il primo livello contiene le porte OR che calcolano i maxtermini e il secondo livello un'unica porta AND che ne fa il prodotto.
\end{itemize}

\begin{esempio}
Sia $f(x,y,z)$ la funzione che vale $1$ solo per $011$, $101$, $110$:
\begin{center}
\begin{tabular}{ccc|c}
\toprule
$x$ & $y$ & $z$ & $f$ \\
\midrule
0 & 0 & 0 & 0 \\
0 & 0 & 1 & 0 \\
0 & 1 & 0 & 0 \\
0 & 1 & 1 & 1 \\
1 & 0 & 0 & 0 \\
1 & 0 & 1 & 1 \\
1 & 1 & 0 & 1 \\
1 & 1 & 1 & 0 \\
\bottomrule
\end{tabular}
\end{center}
La prima forma canonica è $f = x'yz + xy'z + xyz'$. Il circuito ha tre NOT per le variabili negate, tre AND a tre ingressi (una per mintermine) e una OR a tre ingressi.
\end{esempio}

\section{Circuiti speciali: i sommatori}

Con le porte logiche si costruiscono circuiti che eseguono operazioni aritmetiche. Vediamo tre sommatori.

\subsection{Half adder}

Un \emph{half adder} (HA) ha come ingressi due bit $a, b$ e come uscite due bit: $s$ (la somma) e $c$ (il riporto, \emph{carry}).
\begin{center}
\begin{tabular}{cc|cc}
\toprule
$a$ & $b$ & $s$ & $c$ \\
\midrule
0 & 0 & 0 & 0 \\
0 & 1 & 1 & 0 \\
1 & 0 & 1 & 0 \\
1 & 1 & 0 & 1 \\
\bottomrule
\end{tabular}
\end{center}
La colonna $s$ è uno XOR, la colonna $c$ è un AND:
\[
s = a \oplus b, \qquad c = a \cdot b.
\]
Il circuito è quindi formato da una porta XOR e una porta AND con gli stessi ingressi.

\subsection{Full adder}

Un \emph{full adder} (FA) ha tre ingressi, $a$, $b$ e $c_{in}$ (bit e riporto in ingresso), e due uscite, $s$ e $c_{out}$ (somma e riporto in uscita).
\begin{center}
\begin{tabular}{ccc|cc}
\toprule
$a$ & $b$ & $c_{in}$ & $s$ & $c_{out}$ \\
\midrule
0 & 0 & 0 & 0 & 0 \\
0 & 0 & 1 & 1 & 0 \\
0 & 1 & 0 & 1 & 0 \\
0 & 1 & 1 & 0 & 1 \\
1 & 0 & 0 & 1 & 0 \\
1 & 0 & 1 & 0 & 1 \\
1 & 1 & 0 & 0 & 1 \\
1 & 1 & 1 & 1 & 1 \\
\bottomrule
\end{tabular}
\end{center}
Poiché la somma è associativa, $a + b + c_{in} = (a + b) + c_{in}$: si possono usare due half adder, il primo per $a + b$ e il secondo per sommare $c_{in}$ al risultato, e combinare i due riporti con una OR per ottenere $c_{out}$.

\subsection{Sommatore a $n$ bit (ripple-carry)}

Un \emph{sommatore a $n$ bit} ha come ingressi $2n$ bit, $a_{n-1} \dots a_0$ e $b_{n-1} \dots b_0$, e come uscite $n$ bit $s_{n-1} \dots s_0$ più il riporto finale $c$, tali che la sequenza in uscita rappresenti in binario la somma dei due numeri in ingresso. Si realizza collegando $n$ full adder in cascata: il riporto in uscita di ciascuno entra come riporto in ingresso del successivo (da cui il nome \emph{ripple-carry}).

\begin{esempio}
Per sommare due numeri a $3$ bit $X = x_2 x_1 x_0$ e $Y = y_2 y_1 y_0$ servono $3$ full adder in cascata. Ad esempio $100 + 110 = 1010$.
\end{esempio}

\medskip
\noindent\textbf{Riferimenti.} Libro di testo (C. Bolchini et al., \emph{Reti Logiche}, Maggioli): cap.\ 2, sottosezioni 2.6 (porte logiche) e 2.7 (circuiti logici); cap.\ 5, sottosezione 5.2.1 (sommatori ripple-carry). Esercizi 2.24, 2.25, 2.26, 2.27.

\end{document}
